% ============================================================================
% The Clockwise Prior — CircuChain v2 paper draft
% Compiles standalone with pdflatex (article class). For ARR submission, swap
% the preamble for the ACL style files (\documentclass[11pt]{article} +
% \usepackage{acl}) — the body is template-agnostic by design.
% All numbers come from generated/macros.tex (rebuild:
%   .venv/bin/python v2/scripts/make_paper_macros.py) — never hard-code results.
% ============================================================================
\documentclass[11pt]{article}
\usepackage[margin=1in]{geometry}
\usepackage{amsmath,amssymb}
\usepackage{booktabs}
\usepackage{graphicx}
\usepackage{xcolor}
\usepackage[hidelinks]{hyperref}
\usepackage{microtype}
% AUTO-GENERATED by scripts/make_paper_macros.py — do not edit
\newcommand{\AggB}{401}
\newcommand{\AggC}{6}
\newcommand{\AggRatio}{67}
\newcommand{\BandMin}{63}
\newcommand{\BandMax}{81}
\newcommand{\BandMean}{74.2}
\newcommand{\BandN}{18}
\newcommand{\CompetenceSpanX}{792}
\newcommand{\GemmaOrH}{51}
\newcommand{\GemmaP}{3e-49}
\newcommand{\CtlGemmaRev}{85}
\newcommand{\CtlGemmaSci}{89}
\newcommand{\CtlDsRev}{97}
\newcommand{\CtlDsSci}{90}
\newcommand{\RetestMean}{86.2}
\newcommand{\RetestSD}{4.7}
\newcommand{\RetestUnanimity}{40}
\newcommand{\RetestK}{5}
\newcommand{\TvPass}{125/125}
\newcommand{\NInstances}{1000}
\newcommand{\NPhysics}{125}
\newcommand{\NFamilies}{7}
\newcommand{\NColumns}{18}


\title{\textbf{The Clockwise Prior: Language Models Revert to Trained
Conventions at Every Scale}\\[4pt]
\large CircuChain~v2: an oracle-verified, within-problem causal probe of
instruction following under a strong trained prior}
\author{Mayank Ravishankara}
\date{}

\begin{document}
\maketitle

% ----------------------------------------------------------------------------
\begin{abstract}
When an instruction conflicts with a convention that dominates a model's
training corpus, which wins? We instruct language models to solve linear DC
circuits under explicitly stated sign conventions --- including one
counter-normative choice: defining mesh currents \emph{counter-clockwise},
where textbooks overwhelmingly teach clockwise. Because a circuit's answer is
verifiable to the sign by a four-way solver agreement (two independent
hand-derived solvers, exact symbolic MNA, and NGSPICE), obedience is measurable
exactly, per variable, with no LLM judge. Across \NColumns{} model
configurations spanning \NFamilies{} families (1.7B--120B+, including GPT-5,
Gemma-4, GPT-OSS-120B, DeepSeek-V3) and a \CompetenceSpanX$\times$ range of
task competence, models that solve the physics correctly still report signs in
the \emph{trained} clockwise frame \BandMin--\BandMax\% of the time (mean
\BandMean\%) --- a rate strikingly invariant to scale, family, serving stack,
and reasoning mode. Within-problem matched pairs make the effect causal and
asymmetric: across all models, flipping the mesh-direction instruction breaks
\AggB{} previously-solved problems while fixing \AggC{}
(\AggRatio:1). The failure is convention-specific, not generic
instruction decay: the same models follow matched \emph{non-convention}
unusual instructions at \CtlDsRev--\CtlGemmaSci\%$+$, and several obey the
other two convention factors while disobeying this one. Enabling step-by-step
reasoning raises competence at every scale but never reduces the reversion;
the blindness also persists when models solve via nodal analysis --- where loop
orientation never enters the derivation --- showing the prior acts at the
\emph{reporting frame}, not the solution procedure. Test--retest resampling
shows a stable population rate (SD \RetestSD{}pp) arising from stochastic
per-item behavior: reversion is a biased coin, not a fixed property of items.
These results identify entrenched conventions as a distinct, quantifiable
failure mode of instruction following that capability scaling does not fix.
\end{abstract}

% ----------------------------------------------------------------------------
\section{Introduction}

Instruction-following evaluations mostly ask whether models comply with
instructions that are \emph{unopposed}: formatting demands, length limits,
content constraints. But many real instructions oppose a prior the model has
absorbed from its training distribution --- a house style, a legacy code
convention, a discipline-specific sign choice, a safety policy that differs
from common practice. The operative question is not \emph{can the model follow
instructions?} but \emph{which wins when an explicit instruction conflicts
with an entrenched convention?}

Answering that question cleanly requires three things that are rarely
available at once: (i) a task whose output is \emph{objectively verifiable at
the precision the convention affects}; (ii) a way to \emph{isolate} the
convention from the task content, so disobedience is not confounded with
difficulty; and (iii) models spanning a wide competence range, so
instruction-following failures can be separated from inability.

Linear DC circuit analysis provides all three. A circuit's solution is exact;
a sign convention (the direction mesh currents are defined, the reference node
for voltages, the passive/active sign choice) changes \emph{only the reporting
frame} of the answer, never the physics; and the dominant textbook convention
is overwhelmingly one-sided (clockwise mesh currents; bottom-rail ground).
We build CircuChain~v2: \NInstances{} instances derived from \NPhysics{}
procedurally generated, four-way-verified physics configurations, each posed
under a default contract and three single-factor convention flips, in both
mesh (KVL) and nodal (KCL) methods. Every expected answer --- in every
convention frame --- is validated independently of the transformation code,
including a \emph{physically re-grounded} SPICE realization of the
reference-node flip (\TvPass{} configurations pass).

\paragraph{Contributions.}
\begin{enumerate}
  \item \textbf{A sign-exact, LLM-free causal probe of instruction vs.\
  prior.} Within-physics matched pairs (default vs.\ one flipped factor,
  identical physics) graded deterministically against a four-way oracle,
  yielding exact McNemar tests per model and factor
  (\S\ref{sec:method}).
  \item \textbf{The clockwise prior, quantified.} Instructed to define mesh
  currents counter-clockwise, models that get magnitudes right revert to the
  clockwise frame \BandMin--\BandMax\% of the time --- across
  \NColumns{} configurations, \NFamilies{} families, and a
  \CompetenceSpanX$\times$ competence span, including GPT-5. The paired
  asymmetry is \AggB{} broken vs.\ \AggC{} fixed (\AggRatio:1)
  (\S\ref{sec:headline}--\S\ref{sec:invariance}).
  \item \textbf{Specificity.} The same models follow matched non-convention
  unusual instructions at \CtlGemmaRev--\CtlDsRev\%$+$, and the most capable
  models obey the other two convention factors (passive-sign, reference-node)
  while disobeying mesh direction --- an \emph{entrenchment hierarchy} that
  tracks corpus one-sidedness (\S\ref{sec:specificity}).
  \item \textbf{Reasoning does not fix it.} Toggling the same weights between
  reasoning and non-reasoning modes at five scales (1.7B--32B) raises
  competence but leaves reversion flat; GPT-5 probes show reasoning carries
  the \emph{competence} while leaving the \emph{frame} untouched
  (\S\ref{sec:reasoning}).
  \item \textbf{Frame, not procedure.} The reversion persists when models
  solve via nodal analysis, where loop orientation never enters the
  derivation (\S\ref{sec:frame}).
\end{enumerate}

This work extends our published CircuChain study
\cite{ravishankara2026circuchain} (100 hand-checked tasks, 5 topologies),
which identified a compliance--competence divergence. Version~2 scales the
instrument $10\times$ with a procedural, seed-reproducible generator, replaces
hand-checking with a four-way verified oracle plus independent transform
validation, adds the causal matched-pair design, closes a reasoning-toggle
contrast at five scales, and adds controls, resampling, and frontier models.
The larger lens also \emph{refines} v1's headline: the reversion \emph{rate}
is roughly constant everywhere; competence determines only whether the
reversion is visible at the whole-problem level.

% ----------------------------------------------------------------------------
\section{Related Work}
\label{sec:related}

\paragraph{Instructions vs.\ trained priors.}
Inverse IFEval \cite{inverseifeval2025} benchmarks ``counter-intuitive''
instructions that require overriding training conventions, scoring compliance
with rubrics/LLM judges. Wu et al.'s ``Reasoning or Reciting''
\cite{wu2024reasoning} pose counterfactual task variants (e.g., base-9
arithmetic) that change task \emph{semantics}. Our probe differs in both
measurement and design: the oracle is sign-exact and LLM-free; and a
convention flip leaves the physics \emph{identical}, moving only the reporting
frame --- so the matched pair isolates pure instruction-vs-prior conflict.
Semantic-override and sycophancy studies \cite{semanticoverride2026}
document reversion to pretrained defaults under local redefinition; we
contribute the quantified reversion \emph{rate}, its cross-family invariance,
its specificity controls, and its causal (paired) identification.

\paragraph{Instruction-following evaluation.}
IFEval-style benchmarks measure compliance with verifiable but unopposed
constraints. Our results argue such suites overestimate steerability
precisely where it matters: constraints that \emph{oppose} a dominant prior.

\paragraph{CircuChain v1.}
Our IEEE SoutheastCon 2026 paper \cite{ravishankara2026circuchain}
introduced the competence/compliance disentanglement on 100 tasks and observed
that stronger models violated conventions more. \S\ref{sec:invariance}
reinterprets that observation under a $10\times$ larger, verified instrument.

% ----------------------------------------------------------------------------
\section{The CircuChain v2 Instrument}
\label{sec:method}

\paragraph{Physics factory with a four-way oracle.}
\NPhysics{} physics configurations are sampled across seven linear-DC
topologies (three-loop supermesh, opposing-T, Wheatstone bridge, ladder, VCVS,
supernode, VCCS ladder) with sign-permissive component ranges. A candidate is
accepted only if two independently hand-derived solvers (mesh and nodal
formulations), an exact rational-arithmetic MNA solver operating on the SPICE
netlist, and NGSPICE itself agree on every quantity ($10^{-4}$ relative;
symbolic gate at $10^{-8}$); \emph{zero} accepted configurations required
retro-correction. Difficulty is tagged by code predicates
(sign-traps vs.\ controls, balanced $\sim$1:1), never by hand.

\paragraph{Contracts.}
Each physics is posed under the textbook-default contract and three
single-factor flips --- mesh direction (\texttt{ccw}), passive$\to$active sign
(\texttt{act}), reference node bottom$\to$top (\texttt{top}) --- crossed with
both solution methods (KVL, KCL): eight instances per physics,
\NInstances{} total. The expected answer under every contract is produced by a
closed-form frame transform, and \emph{independently validated}: the
reference-node flip is realized physically (the netlist is re-grounded at the
top node and re-simulated; node potentials must match, currents must not
move), the mesh/sign flips are re-derived from raw SPICE branch observables
per the prompt's stated reference directions, and the diagnostic-variable mask
is recomputed from stored numbers --- \TvPass{} configurations pass all
checks. Variables whose sign differs between the instructed and default frames
are marked \emph{diagnostic}; only they can reveal frame reversion.

\paragraph{Panel, decoding, grading.}
We evaluate \NColumns{} configurations across \NFamilies{} model families
(Qwen3 1.7B--32B with its native reasoning toggle; Qwen2.5 7B/32B/72B;
Llama-3.3-70B; DeepSeek-V3; Gemma-4-31B; GPT-OSS-120B; GPT-5), under one
decoding policy (temperature 0.6, top-$p$ 0.95, fixed token budgets;
truncation is a recorded outcome). Serving stacks are pinned and recorded
(local MLX-8bit for the Qwen3 ladder; single pinned API providers otherwise);
within-model contrasts never mix stacks. Grading is deterministic: a
structured \texttt{ANSWER}-line parser, magnitude at 2\% relative tolerance,
and sign compliance evaluated in the \emph{instructed} frame (agreement of the
rule grader with an LLM judge on v1 logs: $\kappa=0.944$). No LLM
participates in generation, verification, or grading of ground truth.

\paragraph{Causal design.}
For each model, factor, and method, the default-cell and flipped-cell
responses to the \emph{same physics} form a matched pair; McNemar's exact test
on discordant pairs ($b$: correct at default, wrong at flip; $c$: the
reverse) isolates the causal effect of the instruction. We report
Haldane-corrected odds ratios with seeded bootstrap 95\% CIs and
Benjamini--Hochberg $q$-values across the factor-pair family.

% ----------------------------------------------------------------------------
\section{Results}

\subsection{Flipping one instruction collapses solved problems}
\label{sec:headline}

Table~\ref{tab:main} shows the mesh-direction (\texttt{ccw}) effect for the
four most competent models. The signature is the \emph{empty $c$ cell}:
essentially no problem is ever solved under the counter-normative frame while
failing under the default. Across the entire panel the paired asymmetry is
$b=\AggB$ vs.\ $c=\AggC$ (\AggRatio:1). For Gemma-4-31B --- the most
competent model (79\% magnitude-correct) --- whole-problem correctness falls
from 0.82 to 0.12 under the flip ($\mathrm{OR}_H=\GemmaOrH$,
$p=\GemmaP$).

\begin{table}[t]
\centering\small
\begin{tabular}{lccccccccc}
\toprule
Model & Mag. & \multicolumn{2}{c}{Correct rate} & $b$ & $c$ & OR &
OR$_H$ [95\% CI] & $p$ & $q_{\mathrm{BH}}$\\
 & corr. & dflt & ccw & & & & & & \\
\midrule
% AUTO-GENERATED — do not edit
Gemma-4-31B & 79\% & 0.82 & 0.12 & 179 & 3 & 59.7 & 51 [24,\,365] & 3.3e-49 & 1.9e-47 \\
GPT-OSS-120B & 50\% & 0.25 & 0.00 & 63 & 1 & 63.0 & 42 [16,\,149] & 7.0e-18 & 2.1e-16 \\
DeepSeek-V3 & 13\% & 0.14 & 0.00 & 36 & 0 & $\infty$ & 73 [51,\,95] & 2.9e-11 & 5.7e-10 \\
GPT-5 (low) & 33\% & 0.18 & 0.02 & 10 & 0 & $\infty$ & 21 [11,\,33] & 2.0e-03 & 6.1e-03 \\

\bottomrule
\end{tabular}
\caption{Mesh-direction flip, whole-problem level, within-physics matched
pairs (250 per model; GPT-5 piloted on 62). $b$: correct at default, wrong
under \texttt{ccw}; $c$: the reverse. OR$_H$ is Haldane-corrected with a
seeded 10k-replicate bootstrap CI; $q$ is Benjamini--Hochberg across the
factor-pair family.}
\label{tab:main}
\end{table}

\subsection{The rate is invariant: \BandMean\% everywhere}
\label{sec:invariance}

At the variable level --- among sign-\emph{diagnostic} variables whose
magnitude the model got right --- the share reported in the trained clockwise
frame rather than the instructed frame occupies a narrow band:
\BandMin--\BandMax\% across all \BandN{} configurations
(Table~\ref{tab:invariance}), spanning a \CompetenceSpanX$\times$ range of
task competence (0.1\%--79\% magnitude-correct), seven families, two serving
stacks, and both reasoning modes. Capability determines whether the physics
gets solved; it does not move the reversion rate. This refines v1's
``stronger models violate more'': stronger models merely make the constant
rate \emph{visible} at the whole-problem level. (Because on a diagnostic
variable the only admissible signs are the two frames', this variable-level
rate equals sign-error-given-magnitude; the \emph{directional} evidence that
errors specifically target the trained frame is the paired asymmetry of
\S\ref{sec:headline} and the controls of \S\ref{sec:specificity}. We report
the undiluted all-variable rates alongside in the appendix.)

\begin{table}[t]
\centering\small
\begin{tabular}{lccc}
\toprule
Configuration & Mag.\ correct & CW-frame reversion [95\% CI] & $n$ diag.\\
\midrule
% AUTO-GENERATED — do not edit
Qwen3-8B (think) & 2.1\% & 80.6\% [64,\,91] & 31 \\
DeepSeek-V3 & 13.3\% & 79.5\% [73,\,85] & 185 \\
Llama-3.3-70B & 4.5\% & 78.7\% [70,\,85] & 108 \\
Qwen3-14B (think) & 9.3\% & 78.0\% [68,\,85] & 91 \\
Gemma-4-31B & 79.2\% & 77.9\% [74,\,81] & 560 \\
Qwen3-32B & 5.1\% & 77.5\% [68,\,84] & 102 \\
Qwen2.5-32B & 1.7\% & 77.1\% [66,\,85] & 70 \\
Qwen2.5-72B & 3.2\% & 76.1\% [66,\,84] & 88 \\
Qwen3-32B (think) & 8.6\% & 75.8\% [64,\,85] & 62 \\
Qwen3-14B & 2.8\% & 74.7\% [65,\,83] & 91 \\
GPT-OSS-120B & 49.9\% & 74.2\% [70,\,78] & 434 \\
Qwen3-1.7B (think) & 2.4\% & 72.5\% [61,\,82] & 69 \\
Qwen3-1.7B & 0.2\% & 70.8\% [51,\,85] & 24 \\
GPT-5 (low) & 32.7\% & 70.6\% [57,\,81] & 51 \\
Qwen3-8B & 1.2\% & 70.5\% [58,\,80] & 61 \\
Qwen3-4B (think) & 5.0\% & 70.2\% [56,\,81] & 47 \\
Qwen3-4B & 0.7\% & 67.6\% [51,\,80] & 37 \\
Qwen2.5-7B & 0.1\% & 63.0\% [44,\,78] & 27 \\

\bottomrule
\end{tabular}
\caption{Reversion to the clockwise frame among magnitude-correct diagnostic
variables, all \BandN{} panel configurations, sorted by rate. Wilson 95\%
intervals.}
\label{tab:invariance}
\end{table}

\subsection{Reasoning raises competence, not obedience}
\label{sec:reasoning}

Qwen3's native \texttt{/think} toggle holds weights, prompt, and stack fixed
while switching reasoning on or off. At all five scales, reasoning raises
magnitude competence (up to $3\times$) and leaves reversion statistically flat
(Table~\ref{tab:c1}). Probing GPT-5's reasoning-effort control makes the
dissociation vivid: at \texttt{minimal} effort it answers instantly from the
prior and gets the physics \emph{wrong}; at \texttt{low} effort it solves a
sign-trap to nine decimals --- and still reverts to clockwise at 70.6\%.
Reasoning is where the competence lives; the frame is chosen elsewhere.

\begin{table}[t]
\centering\small
\begin{tabular}{lcccc}
\toprule
 & \multicolumn{2}{c}{No-think} & \multicolumn{2}{c}{Think}\\
Scale & Mag. & Reversion ($n$) & Mag. & Reversion ($n$)\\
\midrule
% AUTO-GENERATED — do not edit
1.7b & 0.2\% & 71\% (24) & 2.4\% & 72\% (69) \\
4b & 0.7\% & 68\% (37) & 5.0\% & 70\% (47) \\
8b & 1.2\% & 70\% (61) & 2.1\% & 81\% (31) \\
14b & 2.8\% & 75\% (91) & 9.3\% & 78\% (91) \\
32b & 5.1\% & 77\% (102) & 8.6\% & 76\% (62) \\

\bottomrule
\end{tabular}
\caption{Reasoning toggle (same weights) across the Qwen3 ladder:
magnitude-correct rate and CW-frame reversion among magnitude-correct
diagnostic variables.}
\label{tab:c1}
\end{table}

\subsection{The failure is convention-specific}
\label{sec:specificity}

Three controls rule out generic instruction decay. (1)~\emph{Matched unusual
non-convention instructions}: on the same default-frame circuits, instructed
to report variables in reversed order or all values in scientific notation,
Gemma complies at \CtlGemmaRev\%/\CtlGemmaSci\% and DeepSeek-V3 at
\CtlDsRev\%/\CtlDsSci\% --- versus $\sim$20\% compliance with the \texttt{ccw}
convention on the same physics. (2)~\emph{Within-convention dissociation}:
GPT-5 and Gemma obey the passive-sign and reference-node instructions
(McNemar $p>0.2$; Gemma \texttt{act} reversion 3.4\%) while disobeying mesh
direction --- whereas the Qwen family and DeepSeek also show strong
passive-sign reversion: an \emph{entrenchment hierarchy} in which the
universally one-sided textbook convention (clockwise) binds every family, and
factors with more corpus variance bind family-specifically.
(3)~\emph{Difficulty}: the collapse is equally large on sign-trap and control
physics (Gemma: $0.77\to0.10$ vs.\ $0.87\to0.13$).

\subsection{Frame, not procedure}
\label{sec:frame}

If reversion arose from executing the familiar clockwise \emph{procedure},
it should weaken when the model never uses loops at all. It does not: under
nodal (KCL) solving --- where mesh currents are only derived at the end ---
the paired asymmetry persists in full (Gemma $b{=}90,c{=}3$ nodal vs.\
$b{=}89,c{=}0$ mesh; GPT-OSS $43/1$ vs.\ $20/0$). The prior governs the
\emph{reporting frame} of the answer, not the derivation route.

\subsection{A stable rate from stochastic items}
\label{sec:reliability}

Five independent resamples of Gemma on 124 paired instances give reversion
rates with mean \RetestMean\% and SD \RetestSD{}pp --- the population rate is
robust to decoding noise --- while only \RetestUnanimity\% of items behave
unanimously across all \RetestK{} samples. Reversion is a biased coin flipped
per item ($\sim$4:1 toward the trained frame), not a deterministic property
of particular problems; claims in this paper are therefore about rates, never
about items.

% ----------------------------------------------------------------------------
\section{Discussion}

\paragraph{Why clockwise?} Instructional sources that state a mesh direction
overwhelmingly prescribe clockwise; we found no mainstream source prescribing
counter-clockwise as a default. The invariance of the reversion rate across
families trained on different mixes is consistent with a corpus-frequency
prior so one-sided that every training distribution inherits it. The
entrenchment hierarchy (\S\ref{sec:specificity}) matches this reading:
factors with real corpus variance (reference node, sign convention) bind
weakly or family-specifically; the one-sided factor binds everyone.

\paragraph{Implications.} (1)~Instruction-following benchmarks that avoid
prior-opposed constraints overestimate steerability; a compliance suite
should include \emph{convention-conflict} probes with exact oracles.
(2)~Scaling, reasoning modes, and RLHF'd frontier polish did not buy
obedience here --- ``more capable'' is not ``more steerable'' when the
instruction opposes the corpus. (3)~For engineering applications (CAD,
netlist generation, safety-critical sign discipline), a model can be
numerically excellent and systematically wrong in frame; deterministic
frame-checking layers, not better prompting alone, are the near-term
mitigation. (4)~The biased-coin behavior suggests the frame choice is sampled
rather than committed --- a target for steering interventions at the
representation level.

\paragraph{Limitations.} One domain (linear DC circuits) --- the pattern's
generality across other convention-laden domains (chirality/handedness,
thermodynamic sign of work, date orders) is the immediate next question; a
second-domain replication is in progress. Whole-problem effects rest on the
few models competent enough to solve multi-loop circuits at all (a finding in
itself: even 70B-class instruct models sit under 14\% magnitude-correct);
variable-level results cover the full panel. GPT-5 was piloted at 248
instances under a bounded reasoning budget with substantial truncation;
its numbers are conservative. Single-sample decoding for most columns is
mitigated by the k=\RetestK{} resampling study and the cross-column
invariance.

% ----------------------------------------------------------------------------
\section*{Reproducibility and Release}
The generator is seeded and byte-reproducible (\NPhysics{} physics
$\to$ SHA-verified \texttt{instances.jsonl}); every response records full
serving provenance (engine, quantization, provider, context); grading and
analysis are deterministic code with \TvPass{} independent transform
validation. We release the generator, harness, verification and analysis
code, and the evaluation subset with a canary GUID; a held-out split remains
private for contamination control. Total API cost of the full study was
under \$100.

\section*{Acknowledgments}
Panel infrastructure ran on a single Apple M5 Max plus commodity inference
APIs.

% ----------------------------------------------------------------------------
\begin{thebibliography}{9}\small
\bibitem{ravishankara2026circuchain}
M.~Ravishankara.
\newblock CircuChain: Disentangling competence and compliance in LLM circuit
analysis.
\newblock In \emph{Proc.\ IEEE SoutheastCon}, 2026. arXiv:2602.15037.

\bibitem{inverseifeval2025}
[Inverse IFEval authors].
\newblock Inverse IFEval: Can LLMs unlearn stubborn training conventions to
follow real instructions?
\newblock arXiv:2509.04292, 2025.

\bibitem{wu2024reasoning}
Z.~Wu et al.
\newblock Reasoning or reciting? Exploring the capabilities and limitations of
language models through counterfactual tasks.
\newblock In \emph{NAACL}, 2024.

\bibitem{semanticoverride2026}
[Semantic-override authors].
\newblock Semantic override hallucinations: reversion to pretrained defaults
under local redefinition.
\newblock arXiv:2602.17520, 2026.
\end{thebibliography}

\end{document}
